(1) 编码未知量子态 \(|\psi\rangle=\alpha|0\rangle+\beta|1\rangle\) 为 \(|\psi_L\rangle=\alpha|000\rangle+\beta|111\rangle\)。取稳定子 \(S_1=Z_1Z_2,\ S_2=Z_2Z_3\)，列出无错及三个单比特 \(X_i\) 错误的两位综合征，说明测量为何不泄漏 \(\alpha,\beta\)。
(2) 每个物理量子比特独立以概率 \(p\) 遭受 \(X\) 翻转。假设综合征测量与恢复理想，严格求逻辑 \(X_L\) 错误率 \(P_L(p)\)，并说明两位错误为何被错误恢复成三位逻辑翻转。
(3) 求该逻辑bit-flip通道对最大纠缠态的纠缠保真度 \(F_e\)，比较未编码单比特的保真度，精确确定编码改善的 \(p\) 区间，并检查 \(p=0,1/2,1\)。
(4) 改为每个物理量子比特独立概率 \(q\) 发生 \(Z\) 相位翻转。证明所有单 \(Z_i\) 均不触发上述综合征，计算逻辑相位错误率并解释为什么必须换用互补码或更完整的量子纠错码。
